\clearpage
\VerdierPhysicalPageAnchor{0148}%
\begingroup
\small
\setlength{\parskip}{0pt}%
\setlength{\emergencystretch}{1.5em}%
\setlength{\abovedisplayskip}{3pt plus 1pt minus 1pt}%
\setlength{\belowdisplayskip}{3pt plus 1pt minus 1pt}%

\begingroup\itshape
\VerdierSemanticAnchor{ch2-2-3-5-a-ii}
\noindent\hspace{1.2em}ii) 每个态射 $B\longrightarrow X$，其中
$B\in\operatorname{Ob}(\mathcal B)$ 且 $X\in\operatorname{Ob}(\mathcal A)$，
都可嵌入交换图：
\VerdierSemanticAnchor{ch2-2-3-5-a-ii-diagram}
\begin{center}
 \begin{tikzcd}[column sep=4.2em,row sep=2.6em]
 B \arrow[r] \arrow[dr] & X & {B'\in\operatorname{Ob}(\mathcal A\cap\mathcal B)} \\
 & B' \arrow[u] &
 \end{tikzcd}\quad .
\end{center}
\VerdierSemanticAnchor{ch2-2-3-5-a-definition}
$\mathcal D$ 中满足上述两个等价性质的满三角子范畴 $\mathcal A$ 称为
\emph{右 $\mathcal B$-局部化子范畴}。
\VerdierSemanticAnchor{ch2-2-3-5-b}
\par\noindent b) 若 $\mathcal A$ 是右 $\mathcal B$-局部化的，则典范函子
\VerdierSemanticAnchor{ch2-2-3-5-b-functor}
\[
 \mathcal A/(\mathcal A\cap\mathcal B)\longrightarrow\mathcal D/\mathcal B
\]
充分忠实，并且在对象上单射。
\VerdierSemanticAnchor{ch2-2-3-5-c}
\par\noindent c) 下列条件等价：
\VerdierSemanticAnchor{ch2-2-3-5-c-i}
\par\noindent\hspace{1.2em}i) $\mathcal A$ 是右 $\mathcal B$-局部化的。
\VerdierSemanticAnchor{ch2-2-3-5-c-ii}
\par\noindent\hspace{1.2em}ii) 设
$Q\colon\mathcal D/(\mathcal A\cap\mathcal B)\longrightarrow
\mathcal D/\mathcal B$ 为典范函子。子范畴
$\mathcal A/(\mathcal A\cap\mathcal B)\longrightarrow
\mathcal D/(\mathcal A\cap\mathcal B)$ 由 $Q$-右自由对象组成。
\endgroup

\VerdierSemanticAnchor{ch2-2-3-5-proof-a}
a) 中的等价性立即由 TRIII 以及乘法系 $S_{\mathcal D}(\mathcal B)$ 与
$S_{\mathcal A}(\mathcal A\cap\mathcal B)$ 的定义 (2.1.7) 得到。
\VerdierSemanticAnchor{ch2-2-3-5-proof-b-setup}
下面证明 b)。设 $X$ 与 $Y$ 为 $\mathcal A$ 的两个对象（因而也是
$\mathcal A/(\mathcal A\cap\mathcal B)$ 的对象）。
\VerdierSemanticAnchor{ch2-2-3-5-proof-b-injectivity}
先证明同态
\[
 \operatorname{Hom}_{\mathcal A/(\mathcal A\cap\mathcal B)}(X,Y)
 \longrightarrow\operatorname{Hom}_{\mathcal D/\mathcal B}(X,Y)
\]
是单射。
\VerdierSemanticAnchor{ch2-2-3-5-proof-b-quotient-functors}
记 $Q_1\colon\mathcal A\longrightarrow\mathcal A/(\mathcal A\cap\mathcal B)$
与 $Q_2\colon\mathcal D\longrightarrow\mathcal D/\mathcal B$ 为典范商函子。
\VerdierSemanticAnchor{ch2-2-3-5-proof-b-fraction}
任意元素 $m\in\operatorname{Hom}_{\mathcal A/(\mathcal A\cap\mathcal B)}(X,Y)$
可写成 $m=Q_1(s)^{-1}Q_1(f)$，其中
$s\in S_{\mathcal A}(\mathcal A\cap\mathcal B)$ (2.2.2)。它在
$\operatorname{Hom}_{\mathcal D/\mathcal B}(X,Y)$ 中的像是
$Q_2(s)^{-1}Q_2(f)$。
\VerdierSemanticAnchor{ch2-2-3-5-proof-b-zero-criterion}
由 (2.2.4, a))，该像为零，当且仅当存在交换图：
\VerdierSemanticAnchor{ch2-2-3-5-proof-b-zero-diagram}
\begin{center}
 \begin{tikzcd}[column sep=3.7em,row sep=2.5em]
 & {} \arrow[r,"u"] & {} \\
 X \arrow[ur,"f"] && Y \arrow[ul,"s"] \arrow[u,"t"'] &
 {t\in S_{\mathcal D}(\mathcal B)}
 \end{tikzcd}
\end{center}
\endgroup
